% from ega3-1-fr.tex
\oldpage[0\textsubscript{III}]{68}
\emph{Soient $X$ un espace topologique,
$(\mathcal F_k)_{k\in\mathbb N}$ un système projectif de faisceaux de
groupes abéliens sur $X$, et soit
$\mathcal F=\varprojlim_k\mathcal F_k$. On suppose vérifiées les conditions
suivantes :}
\begin{enumerate}
\item[\emph{(i)}]
\emph{Il existe une base $\mathfrak B$ de la topologie de $X$ telle que,
pour tout $U\in\mathfrak B$ et tout $i\geq0$, le système projectif
$(H^i(U,\mathcal F_k))_{k\in\mathbb N}$ vérifie (ML).}
\item[\emph{(ii)}]
\emph{Pour tout $x\in X$ et tout $i>0$, on a}
\[
\varinjlim_U\bigl(\varprojlim_k H^i(U,\mathcal F_k)\bigr)=0
\]
\emph{lorsque $U$ parcourt l'ensemble des voisinages de $x$ appartenant à
$\mathfrak B$.}
\item[\emph{(iii)}]
\emph{Les homomorphismes
$u_{hk}:\mathcal F_k\to\mathcal F_h$ $(h\leq k)$ définissant le système
projectif $(\mathcal F_k)$ sont surjectifs.}
\end{enumerate}
\emph{Dans ces conditions, pour tout $i>0$, l'homomorphisme canonique}
\[
h_i:H^i(X,\mathcal F)\longrightarrow
\varprojlim_k H^i(X,\mathcal F_k)
\]
\emph{est surjectif~; si en outre, pour une valeur de $i$, le système
projectif $(H^{i-1}(X,\mathcal F_k))_{k\in\mathbb N}$ vérifie (ML), $h_i$
est bijectif.}
\end{proposition}

\begin{proof}
\emph{a)} Nous allons d'abord supposer que les $\mathcal F_k$ sont flasques
ainsi que les noyaux $\mathcal N_{hk}$ des $u_{hk}$~; nous allons montrer
alors que la condition (iii) de l'énoncé entraîne que
$H^i(X,\mathcal F)=0$ pour $i>0$. Il suffira de prouver que pour tout ouvert
$U$ de $X$ et tout recouvrement $\mathfrak U$ de $U$ par des ouverts de
$U$, on a $\check H^i(\mathfrak U,\mathcal F)=0$ pour $i>0$. Il en résultera
en effet tout d'abord que pour la cohomologie de Čech, on a
$\check H^i(U,\mathcal F)=0$ pour tout $i>0$, puis (en vertu de
(G, II, 5.9.2) appliqué à l'ensemble de tous les ouverts de $X$) que
$H^i(X,\mathcal F)=0$ pour tout $i>0$. Comme les $\mathcal F_k$ sont
flasques, on a $\check H^i(\mathfrak U,\mathcal F_k)=0$ pour $i>0$
(G, II, 5.2.3)~; considérons pour chaque $k$ le complexe
$\check C^\bullet(\mathfrak U,\mathcal F_k)$ des cochaînes alternées du nerf
du recouvrement $\mathfrak U$ (G, II, 5.1), qui forment évidemment un
système projectif de complexes de groupes abéliens. Montrons que toutes les
applications
$\check C^\bullet(\mathfrak U,\mathcal F_k)\to
\check C^\bullet(\mathfrak U,\mathcal F_h)$ $(h\leq k)$ sont surjectives.
Il suffit évidemment, par définition, de montrer que pour tout ouvert $V$
de $X$, l'application
$\Gamma(V,\mathcal F_k)\to\Gamma(V,\mathcal F_h)$ est surjective~; mais la
suite
\[
0\longrightarrow\mathcal N_{hk}\longrightarrow\mathcal F_k
\longrightarrow\mathcal F_h\longrightarrow0
\]
étant exacte par hypothèse, donne la suite exacte de cohomologie
\[
\Gamma(V,\mathcal F_k)\longrightarrow\Gamma(V,\mathcal F_h)
\longrightarrow H^1(V,\mathcal N_{hk})=0
\]
puisque $\mathcal N_{hk}$ est flasque. Le système projectif
$(\check C^\bullet(\mathfrak U,\mathcal F_k))_{k\in\mathbb N}$ vérifie donc
(ML)~; il en est de même de
$(\check H^i(\mathfrak U,\mathcal F_k))_{k\in\mathbb N}$ pour tout
$i\geq0$, car cela est trivial pour $i>0$, et comme
$\check H^0(\mathfrak U,\mathcal F_k)=\Gamma(U,\mathcal F_k)$
(G, II, 5.2.2), la condition (ML) est aussi remplie pour $i=0$ d'après ce
qui précède. On peut donc appliquer
(\hyperref[0.13.2.3-fr]{13.2.3}), qui montre que
\[
\check H^i(\mathfrak U,\mathcal F)
=\varprojlim_k\check H^i(\mathfrak U,\mathcal F_k)=0
\quad\hbox{pour tout }i>0.
\]

\emph{b)} Passons au cas général, et considérons pour chaque $k\in\mathbb N$,
la résolution canonique
$\mathcal C^\bullet(X,\mathcal F_k)
=(\mathcal C^i(X,\mathcal F_k))_{i\geq0}$ de $\mathcal F_k$ par des
faisceaux flasques (G, II, 4.3). Pour chaque $i\geq0$, il est clair que
$(\mathcal C^i(X,\mathcal F_k))_{k\in\mathbb N}$ est un système projectif
de faisceaux flasques~; montrons qu'il satisfait aux conditions de
\emph{a)}.
